% decay_geiger: a Geiger-Mueller tube (Vol III Ch 2, fig:decay-geiger)
% An ionising particle enters through a thin window and frees a few electrons
% (cPrim) in the gas.  The field of the thin central wire grows as 1/r, so near
% the wire the electrons gain enough energy between collisions to ionise further
% atoms: an avalanche, giving one pulse per particle whatever its energy.
% Electron colour as in decay_antimatter (cPrim); the incoming particle drawn in
% ink because it may be alpha, beta or (indirectly) gamma.
\begin{tikzpicture}[x=1.15cm,y=1.15cm]
\fill[cMoon!60] (0,0) rectangle (6,1.6);
\draw[outline] (0,0) rectangle (6,1.6);
\fill[cTeal!35] (-0.1,0.1) rectangle (0,1.5);
\draw[construction,cAcc] (0.3,0.8) -- (6.7,0.8);
\node[anchor=south west,align=left,small label] at (6.12,0.95) {central wire\\at $+400$\,V};
\node[anchor=south,small label] at (3.0,1.62) {gas-filled metal tube (outer wall earthed)};
\draw[leader,cTeal] (-0.05,0.15) -- (-0.4,-0.3) node[below,small label,text=cTeal] {thin window};
\draw[vector,cInk] (-2.2,1.3) -- (2.3,1.3);
\node[small label,anchor=south east] at (-0.25,1.33) {ionising particle};
\foreach \x in {0.5,0.9,1.3,1.7,2.1}{
  \node[dot,cPrim,minimum size=2.6pt] at (\x,1.3) {};
  \draw[-smallhead,thin line,cPrim] (\x,1.24) -- (\x+0.35,0.9);}
\foreach \x/\y in {2.6/0.92,2.9/0.87,3.15/0.85,3.4/0.84,3.0/0.72,3.3/0.75,3.55/0.77,3.75/0.83,2.8/0.70,3.6/0.90}{
  \node[dot,cPrim,minimum size=2.2pt] at (\x,\y) {};}
\draw[thin line,cPrim,decorate,decoration={brace,amplitude=4pt,mirror}] (2.5,0.6) -- (3.85,0.6)
  node[midway,below=5pt,align=center,small label,text=cInk] {avalanche near the wire:\\one pulse per particle};
\draw[thin line] (6.7,0.8) -- (7.7,0.8) -- (7.7,0.28);
\draw[thin line,fill=white] (7.05,0.28) rectangle (8.35,-0.28);
\node[small label] at (7.7,0) {counter};
\end{tikzpicture}
